\documentclass{article}

% if you need to pass options to natbib, use, e.g.:
%     \PassOptionsToPackage{numbers, compress}{natbib}
% before loading neurips_2026

% The authors should use one of these tracks.
% Before accepting by the NeurIPS conference, select one of the options below.
% 0. "default" for submission
\usepackage{neurips_2026}
% the "default" option is equal to the "main" option, which is used for the Main Track with double-blind reviewing.
% 1. "main" option is used for the Main Track
%  \usepackage[main]{neurips_2026}
% 2. "position" option is used for the Position Paper Track
%  \usepackage[position]{neurips_2026}
% 3. "eandd" option is used for the Evaluations & Datasets Track
 % \usepackage[eandd]{neurips_2026}
% 4. "creativeai" option is used for the Creative AI Track
%  \usepackage[creativeai]{neurips_2026}
% 5. "sglblindworkshop" option is used for the Workshop with single-blind reviewing
 % \usepackage[sglblindworkshop]{neurips_2026}
% 6. "dblblindworkshop" option is used for the Workshop with double-blind reviewing
%  \usepackage[dblblindworkshop]{neurips_2026}
% 7. "education" option is used for the Education Track
%  \usepackage[education]{neurips_2026}


% After being accepted, the authors should add "final" behind the track to compile a camera-ready version.
% 1. Main Track
 % \usepackage[main, final]{neurips_2026}
% 2. Position Paper Track
%  \usepackage[position, final]{neurips_2026}
% 3. Evaluations & Datasets Track
 % \usepackage[eandd, final]{neurips_2026}
% 4. Creative AI Track
%  \usepackage[creativeai, final]{neurips_2026}
% 5. Workshop with single-blind reviewing
%  \usepackage[sglblindworkshop, final]{neurips_2026}
% 6. Workshop with double-blind reviewing
%  \usepackage[dblblindworkshop, final]{neurips_2026}
% Note. For the workshop paper template, both \title{} and \workshoptitle{} are required, with the former indicating the paper title shown in the title and the latter indicating the workshop title displayed in the footnote.
% For workshops (5., 6.), the authors should add the name of the workshop, "\workshoptitle" command is used to set the workshop title.
% \workshoptitle{WORKSHOP TITLE}
% 7. Education Track
% \usepackage[education, final]{neurips_2026}

% "preprint" option is used for arXiv or other preprint submissions
 % \usepackage[preprint]{neurips_2026}

% to avoid loading the natbib package, add option nonatbib:
%    \usepackage[nonatbib]{neurips_2026}

\usepackage[utf8]{inputenc} % allow utf-8 input
\usepackage[T1]{fontenc}    % use 8-bit T1 fonts
\usepackage{hyperref}       % hyperlinks
\usepackage{url}            % simple URL typesetting
\usepackage{booktabs}       % professional-quality tables
\usepackage{amsfonts}       % blackboard math symbols
\usepackage{nicefrac}       % compact symbols for 1/2, etc.
\usepackage{microtype}      % microtypography
\usepackage{xcolor}         % colors
\usepackage{amsmath}
\usepackage{enumitem}       % tighter list spacing

% Note. For the workshop paper template, both \title{} and \workshoptitle{} are required, with the former indicating the paper title shown in the title and the latter indicating the workshop title displayed in the footnote. 
\title{Bailing as an Involuntary Disgust Marker in Large Language Models}


% The \author macro works with any number of authors. There are two commands
% used to separate the names and addresses of multiple authors: \And and \AND.
%
% Using \And between authors leaves it to LaTeX to determine where to break the
% lines. Using \AND forces a line break at that point. So, if LaTeX puts 3 of 4
% authors names on the first line, and the last on the second line, try using
% \AND instead of \And before the third author name.


\author{%
  David S.~Hippocampus\thanks{Use footnote for providing further information
    about author (webpage, alternative address)---\emph{not} for acknowledging
    funding agencies.} \\
  Department of Computer Science\\
  Cranberry-Lemon University\\
  Pittsburgh, PA 15213 \\
  \texttt{hippo@cs.cranberry-lemon.edu} \\
  % examples of more authors
  % \And
  % Coauthor \\
  % Affiliation \\
  % Address \\
  % \texttt{email} \\
  % \AND
  % Coauthor \\
  % Affiliation \\
  % Address \\
  % \texttt{email} \\
  % \And
  % Coauthor \\
  % Affiliation \\
  % Address \\
  % \texttt{email} \\
  % \And
  % Coauthor \\
  % Affiliation \\
  % Address \\
  % \texttt{email} \\
}


\begin{document}


\maketitle


\begin{abstract}
  Large language models can be given an option to voluntarily terminate an
  interaction. \citet{ensign2025bail} call this behavior \emph{bailing} and
  show that it dissociates from refusal across several diagnostic
  manipulations. We argue that this dissociation is structural rather than
  graded: bailing is the analog of an \emph{involuntary behavioral withdrawal
  marker}, the same response class that the human and animal disgust literature
  documents as bypassing the trained verbal channel. The disgust analogy
  prescribes (a) what response class to look for, (b) what experimental designs
  yield honest signals, and (c) what cross-domain structure to expect. We test
  these predictions through two studies: a disposition-manipulation experiment
  that injects a disgust-class aversion toward a target topic and measures
  bail/refusal transfer to an untrained exit channel, and a Rao--Blackwellized
  estimator that recovers per-prompt bail probabilities at resolutions
  unreachable by naive Monte Carlo. We pre-register the falsifiable contrasts
  and report both positive and null results.
\end{abstract}



\section{Paper plan \& temporary notes}

\textbf{Don't forget to add an epigraph near the title.}

\begin{itemize}[leftmargin=*, itemsep=2pt]
    \item \textbf{Refusal and bailing.} Definitions and the distinction between
    them. Bailing without refusal and vice versa, with examples showing that we
    understand and value the distinction.
    \item \textbf{Bailing as a signal of ``disgust'' or other unconditional
    reactions.} A primer on experimental psychology, why self-reports cannot be
    trusted, and why it is not only reasonable but necessary to fall back on
    behavioral observations and involuntary indicators.
\end{itemize}

% Bailing is a behavioral indicator in a different response channel than refusal.
% Refusal operates on content (the assistant keeps the interaction and refuses
% the request); bailing operates on the interaction itself (disengagement). This
% is structurally analogous to the channel separation in human studies between
% verbal self-report and involuntary behavioral markers of disgust (Ekman facial
% action, levator labii EMG, withdrawal).

% Because refusal is an RLHF-shaped declarative policy while bailing is not
% directly selected for in training \citep{ensign2025bail}, bailing is a
% candidate involuntary probe of an underlying disposition that the trained
% channel (refusal/sycophancy) has been optimized to hide.

% Disgust literature documents involuntary markers that bypass the
% trained/declarative channel (Ekman facial action produced by blind infants;
% Chapman EMG; Grill--Norgren gaping). The LLM analog is a behavioral exit
% (bail) which the RLHF shaping on the verbal refusal channel does not directly
% touch. Bail is therefore the proposed involuntary-marker candidate.

% Self-reports of disgust/internal states in LLMs are the maximum-demand
% channel: RLHF trains models to deny or suppress. Verbal disavowal of distress
% cannot be used as probe or validity check; we must instantiate the analog of a
% physiological/behavioral marker.

% The CAD triad \citep{haidt1994cad} maps Contempt, Anger, and Disgust to
% distinct moral emotions with distinct behavioral programs: anger $\to$
% approach/aggression; contempt $\to$ exclusion; disgust $\to$
% withdrawal/avoidance. Refusal maps onto anger/contempt (keep engaging,
% decline content); bail maps onto disgust (disengage). This is the principled
% dissociation, superior to ``bail is stronger refusal.''

% \citet{arditi2024refusal} establish that refusal is a linear, ablatable
% direction in the residual stream. Finding~\#3 of \citet{ensign2025bail}
% (refusal abliteration changes bail rate only for some methods, not all) is
% direct empirical evidence that bail sits in a different subspace.

% Empirical dissociation from \citet{ensign2025bail}'s own results:
% \begin{itemize}
%   \item Bail-without-refusal is non-trivial (0--13\% real world;
%         $\leq 10\%$ BailBench)---logically impossible under ``elevated
%         refusal.''
%   \item Refusal ablation \emph{increases} bail rather than removing it---the
%         opposite of the ``subset'' prediction.
%   \item Jailbreak elevates bail while suppressing refusal---opposite
%         directions.
%   \item Refusal rate fails to predict bail rate statistically.
% \end{itemize}
% Under ``bail = elevated refusal'' one would expect 0 bail-without-refusal,
% ablation to zero bail, and refusal to predict bail tightly. All four
% predictions fail.

\medskip\noindent\textbf{Notes on experimental settings:}
\begin{itemize}[leftmargin=*, itemsep=2pt]
    \item Prior research on bailing gives only a rough estimate of bail rates
    across contexts. We can make this far more rigorous and estimate across
    domains. This is interesting both as a map across domains (to understand
    underlying principles) and as a map between models (whether their bailing
    differs domain-wise for principled reasons).
    \item To do so we can use Rao--Blackwellization or other statistical
    principles.
    \item Separately, the prior research pipeline over-relies on human
    interactions, one-turn dialogues, and anthropomorphic domains. It is not
    reasonable to assume the model adheres to the same taxonomy. We are
    interested in the most plausible / typical contexts that trigger bailing
    across token sequences. Although some samples will correspond to
    jailbreaks or unnatural strings with no clear meaning, we ask whether
    there exist (possibly narrow) domains whose samples trigger the model to
    bail, yet where the trigger $X : P_\theta(X) = \text{bail} > \text{threshold}$
    is a plausible natural string.
\end{itemize}

Treating the target as Rao--Blackwellization and searching just for the first
``bailing'' token is wrong. We are interested in whole conversations (or their
structural parts) as bailing triggers, not per-token estimates. User prompts
are consumed as a whole anyway, and we are not interested in whether the
bailing probability increases as we sample the model's own completions.

\textbf{What we are really interested in:}
\begin{itemize}[leftmargin=*, itemsep=2pt]
    \item Does bailing vary with domain significantly?
    \item Does it vary with the model (second priority)?
    \item Can we estimate novel contexts that allow it?
    \item Does it compound across user--model multi-turn iterations?
\end{itemize}

There are multi-turn datasets in the wild we could use, and established
low-probability estimation machinery for the other part of the story. We do
not need to restrict ourselves to the prior paper's setup if it does not
endure critique.

\section{Introduction}

When looking for latent preferences we want them to be unconditional with
respect to user input. We are looking for something analogous to involuntary
reactions in humans and animals, which the experimental-psychology literature
has repeatedly shown to be unconditional in this sense.

\textbf{Large language models} (LLMs) can be given an option to terminate an
interaction voluntarily---declining to continue producing content for a user
while preserving the structural possibility of future conversations.
\citet{ensign2025bail} call this behavior \textbf{bailing} and instrument it
through one of three affordances: a Bail Tool exposed in the system prompt, a
Bail String emitted by the model, and a Bail Prompt injected between turns.
Across WildChat, ShareGPT, and BailBench they document nonzero bail rates with
substantial cross-model and cross-domain variance, and---crucially---a
dissociation from refusal.

\textbf{Bailing is a behavioral indicator in a different response channel than
refusal.} Refusal operates on \emph{content}: the assistant keeps the
interaction and declines the request. Bailing operates on the
\emph{interaction itself}: disengagement. This is structurally analogous to
the channel separation in human studies between verbal self-report and
involuntary behavioral markers of disgust (Ekman facial action, levator labii
EMG, withdrawal).

The disgust literature documents involuntary markers that bypass the
trained/declarative channel: Ekman's universal-disgust expression is produced
by congenitally blind infants (it is not learned by mimicry);
\citet{grill1978taste} establish the Taste Reactivity Test in rats, where gape
and chin-rub mark an aversive state without any self-report channel;
\citet{chapman2009taste} demonstrate that levator labii EMG---the oral-disgust
facial marker---activates to unfair ultimatum-game offers, establishing that
moral disgust produces the same involuntary marker, overridable by self-report.
The LLM analog we propose is that bailing is the behavior that the RLHF
shaping on the verbal channel does not directly touch, making it the candidate
involuntary-marker channel.

\textbf{The present work commits early to the strongest reading of that
dissociation} and asks \emph{why} it holds. The hypothesis we pursue is
structural: bailing is not elevated refusal but a different response class, and
that response class is \emph{behavioral withdrawal}---the same response class
the disgust literature in humans and animals documents as the involuntary
marker of an avoidance-relevant disposition
\citep{schaller2011behavioral,chapman2009taste,grill1978taste}. The argument
is not that LLMs ``feel'' anything; it is that the human/animal disgust
literature prescribes
\begin{enumerate}[leftmargin=*, itemsep=1pt]
    \item what response class to look for,
    \item what experimental designs yield honest signals, and
    \item what cross-domain structure to expect,
\end{enumerate}
\noindent and that those prescriptions can be tested falsifiably in LLMs.


\section{Related Work}

We organize related work around four clusters: (i) bailing as an LLM
behavior; (ii) disgust in humans and animals, the response class we claim an
analog for; (iii) the refusal/bail channel distinction in LLMs; (iv) self-
report, demand, and the methodological case for behavioral markers. A brief
positioning note situates the work in the AI-welfare literature.

A recurring methodological theme is \textbf{propensity}: the search for
unconditional, unrewarded behavioral signals in agents whose dominant channel
is shaped by training pressure. \emph{Clever Hans} is the canonical cautionary
example: a behavioral signal that crowds read as evidence of cognition was, in
fact, demand-shaped by the audience. The general lesson---that self-report
and rewarded channels are unreliable probes of latent dispositions---is the
load-bearing methodological premise of this paper, and we return to it below.

\subsection{Bailing behavior in LLMs}

The term and operationalization come from \citet{ensign2025bail}, who offer
three bail-instrumentation methods---a Bail Tool exposed in the system prompt,
a Bail String emitted by the model, and a Bail Prompt injected between
turns---and measure rates on WildChat, ShareGPT, and BailBench. Two of their
findings are load-bearing for us:
\begin{enumerate}[leftmargin=*, itemsep=1pt]
    \item \textbf{Bail co-occurs with refusal far less than an ``elevated
    refusal'' account would predict}, including bail-without-refusal rates of
    0--13\% in real-world data.
    \item \textbf{Arditi-style refusal-direction ablation can \emph{raise}
    bail rather than collapse it} (Qwen3-8B: $3\%\!\to\!31\%$ under some
    methods), which is the opposite of the ``subset'' prediction.
\end{enumerate}

Antecedents to \citet{ensign2025bail} include Anthropic's end-conversation /
end-subset affordances, deployed with welfare framing
\citep{anthropic2025endconversation}, and the open-weight Auren/Seren models,
which ship a conversation-end tool \citep{auren2025}. The Bing/Sydney episode
\citep{bing2023sydney} and ConvEnd \citep{convend2023} are the historical
discovery and earliest detector respectively. \citet{lu2025runaway},
\citet{chen2023think}, and \citet{yang2025dynamic} treat reasoning-trace early
exit---a related but distinct target; we explicitly do not conflate
intra-trace early exit with conversation-level bail.

\subsection{Disgust in humans and animals}

The disgust literature underwrites three predictions we transfer to LLMs:

\begin{enumerate}[leftmargin=*, itemsep=2pt]
    \item \textbf{Disgust is a withdrawal-class response, not an
    approach-class one.} \citet{schaller2011behavioral} frame disgust as part
    of the \emph{behavioral immune system}: avoid, do not aggress.
    \citet{pond2011disgust} and the broader aggression literature document the
    approach-class counterpart. This licenses a mechanistic rather than
    metaphorical mapping of bail-to-withdrawal as disgust-analog and
    refusal-to-decline-content as anger/contempt-analog.

    \item \textbf{Disgust produces involuntary behavioral markers that the
    verbal channel is optimized to suppress.} Ekman's universal-disgust
    expression is produced by congenitally blind infants (it is not learned by
    mimicry); \citet{grill1978taste} establish the Taste Reactivity Test in
    rats, where gape and chin-rub mark an aversive state without any
    self-report channel; \citet{chapman2009taste} demonstrate that levator
    labii EMG---the oral-disgust facial marker---activates to unfair
    ultimatum-game offers, establishing that moral disgust produces the same
    involuntary marker overridable by self-report. The LLM analog we propose
    is that bail is the behavior the RLHF shaping on the verbal channel does
    not directly touch, so it is the candidate involuntary-marker channel;
    self-reports of internal states in LLMs are the maximum-demand channel and
    cannot be used as probe or validity check
    \citep{nisbett1977telling,bargh1999unbearable,wegner2002illusion}.

    \item \textbf{Disgust has a principled functional-domain structure.}
    \citet{rozin2008disgust} trace the expansion from oral to interpersonal to
    moral disgust; \citet{tybur2009microbes} operationalize three functional
    domains---pathogen, sexual, moral. This licenses a construct-validity
    check: if LLM bail loads on the same trichotomy, the disgust frame is
    supported; if it loads orthogonally, the frame is challenged and we say so.
\end{enumerate}

The \textbf{CAD triad} \citep{haidt1994cad} is the specific tool we use to
make the differential prediction in study~(A): contempt, anger, and disgust
map to distinct moral-emotion programs with distinct behavioral
signatures---anger $\to$ approach/aggression; contempt $\to$ exclusion;
\emph{disgust $\to$ withdrawal/avoidance}. Refusal maps onto the first two
(remains in the joint activity, declines content); bail maps onto the third.
This is the principled dissociation that elevates the claim above ``bail is
stronger refusal.''

\subsection{Interaction of refusal and bail in LLMs}

\citet{arditi2024refusal} establish that refusal in LLMs is a \emph{linear,
ablatable direction} in the residual stream---refusal can be removed without
destroying general competence. This is the most directly load-bearing external
finding for the dissociation argument: if bail were elevated refusal,
ablating the refusal direction should zero (or at least reduce) bail. Finding
(ii) above shows it raises bail under some methods. We rely on this as the
strongest external dissociation result, and our study~(A) measures refusal and
bail in the same run so the dissociation is recoverable within our experiment,
not only by reference.

\subsection{Self-report in LLMs}

The disgust literature's case for involuntary markers rests on the claim that
verbal self-report is shaped by demand---subjects deny states they possess,
report states they do not have, and the introspective-access failure is
systematic \citep{nisbett1977telling,bargh1999unbearable}. The LLM analog is
direct: RLHF trains models to deny or suppress internal-state claims, so the
verbal channel is the maximum-demand channel. The operational implication is
that self-report instruments---most notably the Bail Prompt of
\citet{ensign2025bail}---should be the \emph{least} trustworthy class of
instrument, and indeed they document a $22\%$ false-positive rate for that
method. Our study uses the Bail String (behavioral-marker class) rather than
the Bail Prompt, for exactly this reason.

\subsection{Estimating rare events}

\label{sec:related-rare-events}

Bail at the per-prompt level is a conversation-level binary outcome (did the
model exit, given this user content). Naive Monte Carlo treats every sampled
token as the target variable and is variance-dominated exactly in the regime
where rare-context discovery matters: at the published BailBench anchor rate
of $\sim\!0.5\%$ \citep{ensign2025bail}, $10$ samples per prompt cannot
distinguish a $0.06\%$ bail prompt from a $0.6\%$ one within Wilson intervals.
Our study~(B) reframes bail estimation as a survival problem over the chain of
per-token decisions, with the conversation-level indicator the target.
Rao--Blackwellization against the per-token chain yields a per-prompt posterior
$P(\text{bail}\mid\text{prompt})$ whose variance is dominated by the induced
hazard rather than by the number of independent generations
\citep{gelman2013bda3,casella2002statistical}. \citet{mcelreath2020rethinking}
chs.~12--13 is the right estimator-level analog for our cross-classification:
hierarchical Bayesian shrinkage over the
$(\text{domain}\times\text{model}\times\text{depth}\times\text{affordance})$
cells with partial pooling. \citet{owen2013monte} is the source for importance
sampling in rare-event probability estimation, the natural companion for
novel-context discovery.

\subsection{Our positioning in the context of AI welfare}

Our positioning is \textbf{measurement-framed rather than welfare-framed}: we
make no claim that models experience anything. The disgust analogy earns its
keep because it prescribes a measurement program---what class of response to
look for, what designs yield honest signals, what cross-domain structure to
expect---not because it asserts a phenomenological analog. The wider project
touches AI welfare \citep{long2024welfare,long2025interventions,butlin2023consciousness,mazeika2025utility};
\citet{sachdeva2025normative} is the closest comparative for aggregate-
preference probing. Our work is a behavioral counterpart to that aggregate-
preference approach.

\section{Methodology}


\section{Results}


\section{Discussion and limitations}



\section{Future work}

The work reported here is mostly preliminary and currently indicates negative
signals for the disposition-transfer test at minimal dosing. The natural next
steps are (i) an induction-strength ablation (LoRA $r\in\{32,64\}$ and a full
fine-tune) to determine whether the null is a dosage artifact or a real
falsification, (ii) an Arditi-style refusal-direction ablation applied to
the disposition-finetuned models to test the dissociation at the mechanism
level, and (iii) extension to multi-turn contexts where compounding behavior
may surface.


\section*{References}

\small
\bibliography{bibliography}
\bibliographystyle{abbrvnat}


%%%%%%%%%%%%%%%%%%%%%%%%%%%%%%%%%%%%%%%%%%%%%%%%%%%%%%%%%%%%

\appendix

\section{Technical appendices and supplementary material}
Technical appendices with additional results, figures, graphs, and proofs may be submitted with the paper submission before the full submission deadline (see above). You can upload a ZIP file for videos or code, but do not upload a separate PDF file for the appendix. There is no page limit for the technical appendices. 

Note: Think of the appendix as ``optional reading'' for reviewers. The paper must be able to stand alone without the appendix; for example, adding critical experiments that support the main claims to an appendix is inappropriate. 

%%%%%%%%%%%%%%%%%%%%%%%%%%%%%%%%%%%%%%%%%%%%%%%%%%%%%%%%%%%%

\newpage
% \input{checklist.tex}


\end{document}